This report reconstructed the paper ``Gaussian Mixture Model Clustering with Incomplete Data'' as a cleaner and more critical technical document. The main algorithmic idea is to optimize the completed data matrix and the GMM parameters jointly, so that imputation is no longer detached from clustering. Once written in a consistent notation, the method becomes easier to understand: the E-step computes responsibilities, the standard GMM M-step updates mixture parameters, and the distinctive data-completion step solves a precision-weighted conditional reconstruction problem for each incomplete sample.

Beyond reproducing the paper, we emphasized what makes the method scientifically interesting. Its empirical advantage comes from aligning missing-value completion with cluster geometry, not from imputation quality in isolation. At the same time, its guarantees remain local, its computational cost is nontrivial, and its practical reliability depends on initialization, covariance regularization, the missingness mechanism, and the choice of the number of clusters.

As a result, this paper should be read not only as a standalone algorithm, but also as a template for a broader design principle in unsupervised learning with incomplete observations: \highlight{latent structure and data completion should be learned together whenever possible}. That principle opens the door to stronger probabilistic, scalable, and missingness-aware extensions in future work.
